(1) 轴向均匀电流密度 $J=I/(\pi a^2)$。安培环路定律 $B_\varphi 2\pi r=\mu_0I_{\rm enc}$ 给 $r<a$ 时 $B_\varphi=\mu_0Ir/(2\pi a^2)$，$a<r<b$ 时 $B_\varphi=\mu_0I/(2\pi r)$；在 $r>b$ 包含回流后总包围电流为零，外部磁场为零。
(2) 能量密度 $u_B=B^2/(2\mu_0)$。圆柱内部每长度 $U'_{\rm in}=\int_0^a[\mu_0I^2r^2/(8\pi^2a^4)]2\pi r\,dr=\mu_0I^2/(16\pi)$；外间隙 $U'_{\rm gap}=\int_a^b[\mu_0I^2/(8\pi^2r^2)]2\pi r\,dr=\mu_0I^2\ln(b/a)/(4\pi)$，两项均非负。
(3) 按定义 $U'=\frac12L'I^2$，得到 $L'=\mu_0/(8\pi)+[\mu_0/(2\pi)]\ln(b/a)$。其中 $\mu_0/(8\pi)$ 来自均匀电流在实体导体内储存的磁能，另一项来自导体间隙；$L'$ 的单位是 $\mathrm{H\,m^{-1}}$。
(4) 当 $b\to a^+$，外间隙趋零而低频均匀载流内电感仍为 $\mu_0/(8\pi)$；当 $b\gg a$，对数间隙贡献逐渐主导。交流高频趋肤效应使电流被挤到圆柱表面，内部磁场与能量减小，理想高频极限的内部电感趋零；不能沿用均匀直流分布公式。
